Se quiere implementar un sistema de votación electrónica en el que se espera que un usuario pueda reemplazar su voto si decide enviarlo nuevamente. Explique qué cambio habría que hacer para conseguir esta propiedad en el sistema de votación basado en firmas de anillo enlazables discutido en clases e implementado en la tarea 4.

\paragraph{Solución.} Para lograr mantener el último voto de cada usuario hacemos lo siguiente. Supo\-nemos que en cada momento tenemos almacenado un conjunto $C$ de votos. Cada vez que llega un nuevo voto $V$, se verifica si este voto fue firmado por el mismo participante que un voto $V'$ que está en $C$. Si es así, entonces se elimina $V'$ de $C$, y se agrega $V$ al conjunto $C$. Si no es así, entonces se agrega $V$ al conjunto $C$ (sin eliminar otros votos).
